% Section 6.2, from lectures/L06/appendix/b_sleep.md.

\section{Sleep}
\label{c6:sec:sleep}

\subsection{What sleeping is}
\label{c6:sec:sleeping}
\code{sleep} is an instruction that stops the processor until something wakes it. What
``something'' can be, and what else stops alongside the processor, depends on the mode you
selected first.

The point is power. An ATmega328P running flat out at 16\,MHz draws on the order of ten
milliamps; in power-down it draws microamps. For anything battery powered that is the difference
between days and years, and it is available to any program whose main loop is mostly waiting.

\textbf{Sleeping is not idling.} A main loop that spins doing nothing is running at full speed and
full current. Replacing that spin with \code{sleep} changes nothing about the program's logic and
everything about its power, which makes it one of the few genuinely free improvements in embedded
work, provided you choose the mode correctly.

\subsection{Setting it up}
\label{c6:sec:smcr}
Two registers and one instruction.

\code{SMCR} holds the mode in \code{SM2:SM0} and the enable in \code{SE}. \code{SE} must be set for
\code{sleep} to do anything at all, and the datasheet's advice is to set it immediately before the
\code{sleep} and clear it immediately after, so that a stray \code{sleep} reached by accident does
nothing.

The layout is one byte with three bits that matter and one that switches it on:

\begin{console}
    bit    7   6   5   4   3   2   1   0
           -   -   -   -  SM2 SM1 SM0  SE
\end{console}

\code{SE} is bit 0 and the mode occupies bits 3 to 1, which means the mode is the encoding
\textbf{shifted left by one}: power-down is \code{SM2:SM0 = 010}, so the byte is \code{0x04}, and
\code{0x05} with \code{SE} set. Power-down written without that shift, \code{010} with \code{SE}
beside it, is \code{0x03}: ADC noise reduction, one mode too shallow, and the program still sleeps,
which is why this is worth writing down rather than deriving each time.

\code{SMCR} is inside the I/O window, so \code{in} and \code{out} reach it, unlike the watchdog's
register: it is data address \code{0x53} and I/O address \code{0x33}, and the name in
\file{m328Pdef.inc} is the second (\secref{c1:sec:iowindow}).

The sequence is: choose the mode, set \code{SE}, execute \code{sleep}, and clear \code{SE} when you
wake. What wakes you runs its interrupt handler first, and then execution continues at the
instruction after the \code{sleep}.

\subsection{What each mode leaves running}
\label{c6:sec:running}

\bookfigure[htbp]{sleep_modes}{Five of the ATmega328P's six sleep modes against four things that might
still be running: idle stops only the CPU, and the deeper modes stop the timers as well, while the
watchdog and pin change interrupts run in all five. Extended standby is
omitted.}{c6:fig:sleep}

The datasheet gives this as a dozen clock domains. Four of them decide whether a given program can
use a given mode, and they are the four in that figure.

\textbf{Idle} stops the CPU and nothing else. Every peripheral keeps running and any interrupt
wakes you. It saves the least and costs nothing in capability, and it is what most programs should
use first.

\textbf{ADC noise reduction} stops the CPU and the I/O clock but leaves the ADC running, so a
conversion happens with the digital noise switched off. It exists for one job and this book does
not do that job; it is in the figure because it sits between idle and power-down in the
\code{SM2:SM0} encoding and you will meet it in the datasheet's table.

\textbf{Power-down} stops nearly everything, including the main clock, and is where the large
savings are. The consequence is the important part: \textbf{a timer cannot wake you from
power-down}, because the timer's clock has stopped too.

\textbf{Power-save} is power-down with Timer2 left running, if Timer2 has been set up to run
asynchronously from a separate crystal. That is the arrangement for a device that must sleep deeply
and still keep time, and it is why Timer2 exists in the form it does. \textbf{An Arduino Uno has no
such crystal}, so on the board this book assumes, power-save wakes from exactly what power-down
wakes from and costs the same to leave; the mode is real and the reason to choose it is not
fitted.

\textbf{Standby} is power-down with the crystal oscillator left running, so waking is faster. It
costs more current than power-down and less than idle. \textbf{Extended standby} is the same trade
applied to power-save, and the figure leaves it out for space: six modes fit in \code{SM2:SM0} and
five of them fit on a page.

So the ladder is six modes deep, from idle through ADC noise reduction, power-down, power-save and
standby to extended standby, and the two that matter for anything in this book are the first and
the third.

\subsection{The two things that always wake you}
\label{c6:sec:wake}
\textbf{A pin change interrupt}, because it is asynchronous: the pin changing is itself the
signal, and no clock is needed to notice it. That is the mechanism behind a device that sleeps
until a button is pressed.

\textbf{The watchdog}, because it runs from its own oscillator and does not care what else has
stopped. Configure it in interrupt mode (\secref{c6:sec:modes}), sleep, and the timeout wakes you.

Those two together are the whole toolkit for a low-power device: \textbf{sleep until something
happens, or until the watchdog says enough time has passed.} A program that wants to do something
every eight seconds and nothing in between can do exactly that, and spend almost all of its life
drawing microamps.

\textbf{And this is the thing to get right.} A program that sleeps in power-down and expects Timer1
to wake it in a second does not wake up. Nothing reports an error; the device simply sits there,
drawing microamps, for ever. It is one of the more memorable ways to lose an afternoon, and it is
entirely predictable from Figure~\ref{c6:fig:sleep}.

\subsection{What sleep does not do}
\label{c6:sec:sleepnot}
\textbf{It does not stop time.} Or rather, it stops the timers, which is exactly the problem: a
program that sleeps for a while and then reasons about elapsed time from a timer count is
reasoning from a counter that was not running.

\textbf{It does not save anything if you wake constantly.} Waking has a cost in cycles, and in the
deeper modes it is dominated by waiting for the oscillator to stabilise: a millisecond, at
16\,MHz, which \secref{c6:sec:restart} puts a number on. A program that sleeps and wakes a thousand
times a second in power-down spends all of its time starting an oscillator, and may use more
energy than one that simply idles.

\textbf{It does not survive being interrupted by your own peripherals.} A UART receiving in the
background, a timer you forgot to stop: any enabled interrupt that is still capable of firing will
wake you, and a device that mysteriously refuses to stay asleep is usually being woken by
something it configured and forgot.

\textbf{And it is not modelled here.} This book does not test sleep, because the interaction
between sleep modes, clock domains and wake sources is exactly the kind of thing a simulator
approximates and a datasheet specifies. The material is here because you will need it and because
the power-down trap is worth meeting on paper first; the checking is yours to do on real hardware.

\subsection{What it costs to come back}
\label{c6:sec:restart}
Waking is not free in any mode, and in the deep ones the cost is almost entirely the oscillator.

\begin{narrowtable}{@{}lcrr@{}}
  \thead{Mode} & \thead{\code{SM2:SM0}} & \thead{Oscillator restart} & \thead{At 16\,MHz} \\
  \midrule
  Idle & \code{000} & 0 cycles & 0\,µs \\
  ADC noise reduction & \code{001} & 0 cycles & 0\,µs \\
  Power-down & \code{010} & 16384 cycles & 1024\,µs \\
  Power-save & \code{011} & 16384 cycles & 1024\,µs \\
  Standby & \code{110} & 6 cycles & 0.375\,µs \\
  Extended standby & \code{111} & 6 cycles & 0.375\,µs \\
\end{narrowtable}

The cycle counts are pinned in \file{atmega328p.hpp} as \code{SleepModes}. Work the right-hand
column out before you read it: one cycle is 62.5 nanoseconds, so it is a single multiplication, and
the answer is more surprising than the arithmetic.

\textbf{The 16384 is a fuse setting, not a property of the part.} It is the \code{16K CK} start-up
time the \code{SUT} and \code{CKSEL} fuses select for a crystal, chosen so that the oscillator has
certainly settled before any instruction executes. A part running from its internal RC oscillator
pays six cycles instead, because there is nothing mechanical to settle. So the largest single term
in a power-down wake-up is a decision somebody made with a fuse programmer, and it is \textbf{not
in the source code}.

The two standby modes exist precisely to buy that term back: they are power-down and power-save
with the oscillator left running, so waking costs six cycles rather than 16384. What they cost in
return is the oscillator's current for the whole time the part is asleep, which is most of the
time, which was the point of sleeping. That is the trade in its clearest form, \textbf{standby
spending current all night to save a millisecond in the morning}, and it is a genuine engineering
decision rather than a right answer.

\textbf{Waking up is not the same event as running again}, and the gap between them is not in this
table. What wakes you is an interrupt, so the first thing that runs is a handler, and your program
continues at the instruction after \code{sleep} only once that handler has returned. On a bare
program that is the handler's own cost and nothing else. Under a scheduler it is the handler, plus
whichever task the scheduler then decides should run, which need not be the one that went to
sleep. That is the whole of the difference between a wake-up \emph{latency} and a wake-up
\emph{cost}, and the kernel course this book leads to spends a lecture on it.
